\documentclass[11pt]{article}

\usepackage[T1]{fontenc}
\usepackage[utf8]{inputenc}
\usepackage[english]{babel}
\usepackage{amsmath,amssymb,amsthm,mathtools}
\usepackage[a4paper,margin=2.7cm]{geometry}
\usepackage{microtype}

\newcommand{\R}{\mathbb R}
\newcommand{\Sym}{\mathbb S}
\newcommand{\dd}{\,\mathrm d}
\newcommand{\cP}{\mathcal P}
\newcommand{\aff}{\operatorname{aff}}
\newcommand{\co}{\operatorname{co}}
\newcommand{\Ran}{\operatorname{Ran}}
\newcommand{\tr}{\operatorname{tr}}

\newtheorem{definition}{Definition}
\newtheorem{proposition}{Proposition}
\newtheorem{theorem}{Theorem}
\newtheorem{corollary}{Corollary}
\theoremstyle{remark}
\newtheorem{remark}{Remark}
\newtheorem{example}{Example}

\title{Moment Closure for Wasserstein Gradient Flows}
\author{}
\date{}

\begin{document}

\maketitle

\begin{abstract}
We characterize when the Wasserstein gradient flow of every functional
\(
F_f(\mu)=f(\int h\,\dd\mu)
\)
induces an autonomous dynamics on the selected moment
\(
m_h(\mu)=\int h\,\dd\mu.
\)
The exact condition is that the Gram matrix
\(
Dh(x)Dh(x)^\top
\)
depend affinely on \(h(x)\).  We derive the resulting finite-dimensional
gradient system, classify scalar and vector-valued quadratic observables,
and recover the closed dynamics of the second-moment matrix.  We also give
genuinely non-polynomial examples based on distance functions and block
radii, and show that polynomial observables satisfying the universal
closure condition necessarily have degree at most two.
\end{abstract}

\section{Setup and notion of closure}

Let \(\Omega\subset\R^n\) be open and let
\[
  h=(h_1,\ldots,h_d):\Omega\longrightarrow\R^d
\]
be \(C^2\).  We work with an admissible class of probability measures
\(\cP_h(\Omega)\) for which all the integrals below are finite.  We assume
that this class contains the finitely supported probability measures.  This
last assumption is used only in the necessity part of the closure theorem.

Define the selected moment
\[
  m_h(\mu):=\int_\Omega h(x)\,\dd\mu(x)\in\R^d
\]
and, for \(f:\R^d\to\R\), the functional
\[
  F_f(\mu):=f(m_h(\mu)).
\]
The descending Wasserstein gradient flow of \(F_f\) is formally
\[
  \partial_t\mu_t+\nabla\cdot(\mu_t v_t)=0,
  \qquad
  v_t(x)
  =
  -\nabla_x\frac{\delta F_f}{\delta\mu}(\mu_t)(x).
\]
All statements below can either be read for smooth solutions, or under
standard assumptions ensuring that the displayed differentiation formulas
hold for weak solutions.

\begin{definition}[Universal closure on a selected moment]
The observable \(h\) has the \emph{universal moment-closure property} if,
for every smooth \(f\), there exists a vector field
\[
  b_f:m_h(\cP_h(\Omega))\longrightarrow\R^d
\]
such that every Wasserstein gradient flow of \(F_f=f\circ m_h\) satisfies
\[
  \frac{\dd}{\dd t}m_h(\mu_t)=b_f(m_h(\mu_t)).
\]
Thus the evolution of \(m_h(\mu_t)\) is independent of all information about
\(\mu_t\) not contained in \(m_h(\mu_t)\).
\end{definition}

\begin{remark}[Closure for a family of objectives]
Universal closure is stronger than closure for one fixed \(f\).  For a
fixed objective, only the product of the mobility matrix with
\(\nabla f(m)\) must be determined by \(m\).  Requiring closure for every
\(f\) forces the entire mobility matrix to be determined by \(m\).
\end{remark}

\section{Exact affine Gram criterion}

The relevant object is the carré-du-champ, or Gram matrix, of the selected
statistics:
\[
  \Gamma_h(x):=Dh(x)Dh(x)^\top\in\Sym^d_+,
  \qquad
  (\Gamma_h(x))_{ij}
  =
  \langle\nabla h_i(x),\nabla h_j(x)\rangle.
\]

\begin{proposition}[Evolution of the selected moment]
\label{prop:moment-evolution}
Along the Wasserstein gradient flow of \(F_f=f\circ m_h\),
\[
  \frac{\dd}{\dd t}m_h(\mu_t)
  =
  -G_h(\mu_t)\nabla f(m_h(\mu_t)),
  \qquad
  G_h(\mu):=\int_\Omega\Gamma_h(x)\,\dd\mu(x).
\]
\end{proposition}

\begin{proof}
The first variation of the \(i\)-th component of \(m_h\) is \(h_i\).
Consequently,
\[
  \frac{\delta F_f}{\delta\mu}(\mu)(x)
  =
  \langle\nabla f(m_h(\mu)),h(x)\rangle
\]
and
\[
  v(x)=-Dh(x)^\top\nabla f(m_h(\mu)).
\]
Using the continuity equation, for each \(i\),
\[
  \begin{aligned}
  \frac{\dd}{\dd t}m_{h,i}(\mu_t)
  &=
  \int_\Omega\langle\nabla h_i(x),v_t(x)\rangle\,\dd\mu_t(x)\\
  &=
  -\sum_{j=1}^d
  \partial_j f(m_h(\mu_t))
  \int_\Omega
  \langle\nabla h_i(x),\nabla h_j(x)\rangle\,\dd\mu_t(x).
  \end{aligned}
\]
Collecting the components gives the result.
\end{proof}

\begin{theorem}[Universal moment-closure criterion]
\label{thm:affine-gram}
The following statements are equivalent.
\begin{enumerate}
  \item The observable \(h\) has the universal moment-closure property.

  \item Whenever \(m_h(\mu)=m_h(\nu)\), one has
  \[
    G_h(\mu)=G_h(\nu).
  \]

  \item There exists an affine map
  \[
    A:\aff(h(\Omega))\longrightarrow\Sym^d
  \]
  such that
  \[
    \Gamma_h(x)=A(h(x))
    \qquad\text{for every }x\in\Omega.
  \]
\end{enumerate}
Equivalently, there are constant symmetric matrices
\(C_0,C_1,\ldots,C_d\in\Sym^d\) such that
\[
  \boxed{
  Dh(x)Dh(x)^\top
  =
  C_0+\sum_{k=1}^d h_k(x)C_k.
  }
\]
The matrices are unique only modulo affine relations satisfied by
\(h(\Omega)\).  When these conditions hold, the closed equation is
\[
  \boxed{
  \frac{\dd}{\dd t}m(t)
  =
  -A(m(t))\nabla f(m(t)),
  \qquad
  A(m)=C_0+\sum_{k=1}^d m_kC_k.
  }
\]
Moreover, \(A(m)\) is positive semidefinite for every realizable moment
\(m\in m_h(\cP_h(\Omega))\).
\end{theorem}

\begin{proof}
Suppose first that universal closure holds, and let \(\mu,\nu\) have the
same selected moment \(m\).  Apply Proposition
\ref{prop:moment-evolution} to the linear objective
\(f_a(y)=\langle a,y\rangle\).  Since the closed velocity at \(m\) is the
same for \(\mu\) and \(\nu\),
\[
  G_h(\mu)a=G_h(\nu)a
  \qquad\text{for every }a\in\R^d.
\]
Hence \(G_h(\mu)=G_h(\nu)\), proving \(1\Rightarrow2\).

Assume \(2\).  For a finite convex combination of points in \(h(\Omega)\),
define
\[
  A\left(\sum_{\alpha=1}^N p_\alpha h(x_\alpha)\right)
  :=
  \sum_{\alpha=1}^N p_\alpha\Gamma_h(x_\alpha).
\]
Statement \(2\), applied to finitely supported measures, shows that this
definition is independent of the chosen representation.  It also shows
that \(A\) preserves finite convex combinations.  Therefore \(A\) is affine
on \(\co(h(\Omega))\), and it extends uniquely to the affine hull of
\(h(\Omega)\).  Applying the definition to a Dirac mass gives
\(\Gamma_h(x)=A(h(x))\), proving \(2\Rightarrow3\).

Finally, if \(\Gamma_h=A\circ h\) with \(A\) affine, then
\[
  G_h(\mu)
  =
  \int_\Omega A(h(x))\,\dd\mu(x)
  =
  A\left(\int_\Omega h(x)\,\dd\mu(x)\right)
  =
  A(m_h(\mu)).
\]
Proposition \ref{prop:moment-evolution} gives the closed equation and proves
\(3\Rightarrow1\).  Positive semidefiniteness follows from
\[
  A(m_h(\mu))
  =
  \int_\Omega Dh(x)Dh(x)^\top\,\dd\mu(x)\succeq0.
\]
\end{proof}

\begin{remark}[A finite-dimensional metric]
On a regular region where \(A(m)\) is positive definite, the closed equation
is the gradient flow of \(f\) for the cometric \(A(m)\).  Degeneracy of
\(A(m)\) records differential relations among the selected statistics.
\end{remark}

\section{Scalar observables}

\begin{corollary}[Scalar closure criterion]
\label{cor:scalar}
For \(h:\Omega\to\R\), universal closure holds if and only if there are
constants \(\alpha,\beta\in\R\) such that
\[
  \boxed{
  |\nabla h(x)|^2=\alpha+\beta h(x).
  }
\]
The closed equation is
\[
  \dot m(t)
  =
  -\bigl(\alpha+\beta m(t)\bigr)f'(m(t)).
\]
\end{corollary}

\begin{proof}
For \(d=1\), every affine map on \(\aff(h(\Omega))\subset\R\) has the form
\(A(s)=\alpha+\beta s\).  The result follows from Theorem
\ref{thm:affine-gram}.
\end{proof}

\begin{proposition}[Eikonal normal form on regular regions]
\label{prop:eikonal}
Let \(h\in C^2(\Omega)\) be scalar-valued.  On every connected open set
\[
  U\subset\{x\in\Omega:\nabla h(x)\neq0\},
\]
the following statements are equivalent:
\begin{enumerate}
  \item \(h|_U\) has the universal moment-closure property;
  \item there exist \(\rho\in C^2(U)\) and constants
  \(c,\alpha,\lambda\in\R\) such that
  \[
    |\nabla\rho(x)|^2=1
    \qquad\text{and}\qquad
    h(x)
    =
    c+\alpha\rho(x)+\frac{\lambda}{2}\rho(x)^2.
  \]
\end{enumerate}
For such a representation,
\[
  |\nabla h(x)|^2
  =
  \alpha^2-2\lambda c+2\lambda h(x),
\]
and the closed dynamics is
\[
  \boxed{
  \dot m
  =
  -\bigl(\alpha^2-2\lambda c+2\lambda m\bigr)f'(m).
  }
\]
\end{proposition}

\begin{proof}
Suppose first that \(h|_U\) has universal closure.  By Corollary
\ref{cor:scalar}, there are constants \(a,b\in\R\) such that
\[
  |\nabla h|^2=a+bh.
\]
The right-hand side is strictly positive on \(U\).  Fix \(s_0\in h(U)\)
and define
\[
  \rho(x)
  :=
  \int_{s_0}^{h(x)}
  \frac{\dd s}{\sqrt{a+bs}}.
\]
Then
\[
  \nabla\rho(x)
  =
  \frac{\nabla h(x)}{\sqrt{a+bh(x)}},
  \qquad
  |\nabla\rho(x)|^2=1.
\]
If \(b=0\), then \(a>0\) and
\[
  h=s_0+\sqrt a\,\rho.
\]
If \(b\neq0\), direct integration gives
\[
  \rho
  =
  \frac{2}{b}
  \left(
    \sqrt{a+bh}-\sqrt{a+bs_0}
  \right).
\]
Writing \(\alpha=\sqrt{a+bs_0}\) and \(\lambda=b/2\), this rearranges as
\[
  h=s_0+\alpha\rho+\frac{\lambda}{2}\rho^2.
\]

Conversely, suppose that \(h=c+\alpha\rho+\lambda\rho^2/2\) and that
\(|\nabla\rho|=1\).  Then
\[
  |\nabla h|^2
  =
  (\alpha+\lambda\rho)^2
  =
  \alpha^2-2\lambda c+2\lambda h.
\]
Corollary \ref{cor:scalar} gives universal closure and the displayed
equation.
\end{proof}

\begin{remark}[Critical sets]
The regularity restriction in Proposition \ref{prop:eikonal} is essential
for a smooth eikonal coordinate.  For example,
\[
  h(x)=\frac12|x|^2
\]
is universally closed on all of \(\R^n\), but its natural eikonal coordinate
\(\rho(x)=|x|\) is not \(C^2\) at the origin.  Thus the proposition is a
complete local classification on regular strata.  A global classification
across critical sets requires either additional geometric hypotheses or
nonsmooth distance coordinates.
\end{remark}

\begin{proposition}[Classification of scalar quadratic observables]
\label{prop:scalar-quadratic}
Let
\[
  h(x)
  =
  \langle a,x\rangle+\frac12\langle Bx,x\rangle,
  \qquad
  a\in\R^n,\quad B\in\Sym^n.
\]
If \(B=0\), universal closure always holds and
\[
  \dot m=-|a|^2f'(m).
\]
If \(B\neq0\), universal closure holds if and only if
\[
  B=\lambda P,
  \qquad
  P^2=P=P^\top,
  \qquad
  a\in\Ran P,
\]
for some nonzero \(\lambda\in\R\).  In that case,
\[
  |\nabla h|^2=|a|^2+2\lambda h
\]
and
\[
  \dot m
  =
  -\bigl(|a|^2+2\lambda m\bigr)f'(m).
\]
Thus all nonzero eigenvalues of \(B\) must coincide.
\end{proposition}

\begin{proof}
By Corollary \ref{cor:scalar}, closure is equivalent to
\[
  |a+Bx|^2=\alpha+\beta h(x).
\]
Identifying the constant, linear, and quadratic terms gives
\[
  \alpha=|a|^2,
  \qquad
  2Ba=\beta a,
  \qquad
  B^2=\frac{\beta}{2}B.
\]
If \(B\neq0\), set \(\lambda=\beta/2\).  Since \(B\) is symmetric,
\(B^2=\lambda B\) means that its eigenvalues belong to
\(\{0,\lambda\}\), with \(\lambda\neq0\).  Hence \(B=\lambda P\) for an
orthogonal projection \(P\).  The linear identity becomes \(Pa=a\).
Conversely, these identities directly give
\(|\nabla h|^2=|a|^2+2\lambda h\).
\end{proof}

\section{Vector-valued quadratic observables}

\begin{proposition}[Quadratic Jordan criterion]
\label{prop:vector-quadratic}
Consider
\[
  h_i(x)
  =
  \langle a_i,x\rangle+\frac12\langle B_i x,x\rangle,
  \qquad
  a_i\in\R^n,\quad B_i\in\Sym^n,
  \qquad i=1,\ldots,d.
\]
Universal closure holds if and only if, for every \(i,j\), there are
coefficients \(c_{ij}^1,\ldots,c_{ij}^d\) such that
\[
  \boxed{
  B_iB_j+B_jB_i
  =
  \sum_{k=1}^d c_{ij}^kB_k,
  }
\]
and
\[
  \boxed{
  B_i a_j+B_j a_i
  =
  \sum_{k=1}^d c_{ij}^k a_k.
  }
\]
The affine mobility is then
\[
  A_{ij}(m)
  =
  \langle a_i,a_j\rangle
  \sum_{k=1}^d c_{ij}^k m_k,
\]
and the closed dynamics is
\[
  \dot m=-A(m)\nabla f(m).
\]
In particular, \(\operatorname{span}\{B_1,\ldots,B_d\}\) must be closed
under the Jordan product
\[
  B_i\circ B_j:=\frac12(B_iB_j+B_jB_i).
\]
Jordan closure of the quadratic parts is necessary, while the second
displayed identity is the additional compatibility condition for the
linear parts.
\end{proposition}

\begin{proof}
One has
\[
  \nabla h_i(x)=a_i+B_ix
\]
and therefore
\[
  \begin{aligned}
  \langle\nabla h_i(x),\nabla h_j(x)\rangle
  ={}&
  \langle a_i,a_j\rangle
  \langle B_i a_j+B_j a_i,x\rangle\\
  &+
  \frac12
  \langle(B_iB_j+B_jB_i)x,x\rangle.
  \end{aligned}
\]
By Theorem \ref{thm:affine-gram}, this expression must equal
\[
  c_{ij}^0+\sum_{k=1}^d c_{ij}^k h_k(x).
\]
Identifying constant, linear, and quadratic coefficients gives
\[
  c_{ij}^0=\langle a_i,a_j\rangle
\]
and the two claimed identities.  The converse follows by substituting
these identities back into the Gram matrix.
\end{proof}

\section{The closed second-moment dynamics}

\begin{proposition}[Second-moment matrix]
\label{prop:second-moment}
Let
\[
  M(\mu):=\int_{\R^n}xx^\top\,\dd\mu(x)\in\Sym^n_+
\]
and \(F_f(\mu)=f(M(\mu))\), where
\(f:\Sym^n\to\R\) is differentiated using the Frobenius inner product.
Then \(M(\mu_t)\) satisfies the autonomous equation
\[
  \boxed{
  \dot M
  =
  -2\bigl(\nabla f(M)M+M\nabla f(M)\bigr).
  }
\]
\end{proposition}

\begin{proof}
For \(U\in\Sym^n\), the scalar component
\[
  h_U(x):=\langle U,xx^\top\rangle_{\mathrm F}=x^\top Ux
\]
satisfies \(\nabla h_U(x)=2Ux\).  Hence, for \(U,V\in\Sym^n\),
\[
  \begin{aligned}
  \langle\nabla h_U(x),\nabla h_V(x)\rangle
  &=
  4x^\top UVx\\
  &=
  2x^\top(UV+VU)x\\
  &=
  \langle2(UV+VU),xx^\top\rangle_{\mathrm F}.
  \end{aligned}
\]
The Gram form is linear in \(xx^\top\), so Theorem
\ref{thm:affine-gram} already gives closure.

For the explicit equation, set \(G=\nabla f(M)\in\Sym^n\).  Then
\[
  \frac{\delta F_f}{\delta\mu}(x)=x^\top Gx,
  \qquad
  v(x)=-2Gx.
\]
Differentiating \(M\) along the continuity equation gives
\[
  \begin{aligned}
  \dot M
  &=
  \int_{\R^n}\bigl(v(x)x^\top+xv(x)^\top\bigr)\,\dd\mu(x)\\
  &=
  -2(GM+MG),
  \end{aligned}
\]
as claimed.
\end{proof}

\begin{remark}[Mean and second moment together]
The augmented statistic
\[
  h(x)=(x,xx^\top)
\]
also satisfies the affine Gram criterion.  Thus objectives depending jointly
on the mean and the uncentered second moment induce a closed dynamics on
that pair.  The covariance is then recovered as
\(M-\bar x\bar x^\top\).
\end{remark}

\section{Genuinely non-quadratic closed moments}

Non-quadratic examples exist, but on Euclidean space they are naturally
non-polynomial and are often local because distance functions develop
singularities outside a tubular neighborhood.

\begin{example}[Distance and signed-distance moments]
Proposition \ref{prop:eikonal} gives the following non-quadratic examples.
\begin{itemize}
  \item On \(\R^n\setminus\{x_0\}\), take
  \(\rho(x)=|x-x_0|\).  In particular,
  \[
    m(\mu)=\int|x-x_0|\,\dd\mu(x)
  \]
  satisfies \(\dot m=-f'(m)\).

  \item More generally, the distance to an affine subspace has unit gradient
  away from the subspace and yields a closed radial moment.

  \item If \(\Sigma\) is a \(C^2\) hypersurface, its signed distance
  \(\rho_\Sigma\) is \(C^2\) and satisfies
  \(|\nabla\rho_\Sigma|=1\) on a sufficiently small tubular neighborhood.
  Hence both
  \[
    \int\rho_\Sigma(x)\,\dd\mu(x)
    \quad\text{and}\quad
    \frac12\int\rho_\Sigma(x)^2\,\dd\mu(x)
  \]
  have closed dynamics as long as the flow remains in that neighborhood.
  For a curved hypersurface these are genuinely non-quadratic functions of
  the Euclidean coordinates.

  \item Any affine-quadratic function of a signed distance,
  \[
    \int
    \left(
      \alpha\rho_\Sigma(x)
      +\frac{\lambda}{2}\rho_\Sigma(x)^2
    \right)\dd\mu(x),
  \]
  is closed by the same proposition.
\end{itemize}
\end{example}

\begin{proposition}[Orthogonal block-radius moments]
\label{prop:block-radius}
Let
\[
  \R^n=E_0\oplus E_1\oplus\cdots\oplus E_d
\]
be an orthogonal decomposition, and let \(P_i\) be the orthogonal projection
onto \(E_i\).  Fix \(x_0\in\R^n\) and work on
\[
  \Omega
  :=
  \{x:P_i(x-x_0)\neq0,\ i=1,\ldots,d\}.
\]
Set
\[
  \rho_i(x):=|P_i(x-x_0)|,
  \qquad
  h_i(x)
  :=
  \alpha_i\rho_i(x)+\frac{\lambda_i}{2}\rho_i(x)^2.
\]
Then
\[
  \langle\nabla h_i,\nabla h_j\rangle
  =
  \delta_{ij}\bigl(\alpha_i^2+2\lambda_i h_i\bigr),
\]
so the vector \(m=(\int h_i\,\dd\mu)_{i=1}^d\) has the closed dynamics
\[
  \boxed{
  \dot m_i
  =
  -\bigl(\alpha_i^2+2\lambda_i m_i\bigr)\partial_i f(m),
  \qquad i=1,\ldots,d.
  }
\]
\end{proposition}

\begin{proof}
For every \(i\),
\[
  \nabla\rho_i(x)
  =
  \frac{P_i(x-x_0)}{|P_i(x-x_0)|}\in E_i.
\]
The orthogonality of the blocks gives
\[
  \langle\nabla\rho_i,\nabla\rho_j\rangle=\delta_{ij}.
\]
The same computation as in Proposition \ref{prop:eikonal} then yields
\[
  \langle\nabla h_i,\nabla h_j\rangle
  =
  \delta_{ij}(\alpha_i+\lambda_i\rho_i)^2
  =
  \delta_{ij}(\alpha_i^2+2\lambda_i h_i).
\]
Theorem \ref{thm:affine-gram} gives the displayed dynamics.
\end{proof}

\begin{example}[A short catalogue]
The preceding results provide several explicit non-quadratic closed
statistics:
\begin{enumerate}
  \item the mean radius \(\int|x|\,\dd\mu(x)\) on
  \(\R^n\setminus\{0\}\);
  \item simultaneous block radii
  \(\bigl(\int|P_i x|\,\dd\mu(x)\bigr)_{i=1}^d\);
  \item signed-distance moments around a curved hypersurface;
  \item affine-quadratic functions of these distances;
  \item mixtures of radial and quadratic-radial coordinates on mutually
  orthogonal blocks.
\end{enumerate}
All of these examples are exact closures for every outer objective \(f\),
not closures tied to one particular choice of \(f\).
\end{example}

\begin{proposition}[No higher-degree polynomial examples]
\label{prop:no-higher-polynomial}
Assume that \(\Omega\) contains a nonempty open subset of \(\R^n\), that
each component \(h_i\) is a polynomial, and that \(h\) has the universal
moment-closure property.  Then every \(h_i\) has degree at most two.  Thus
the non-quadratic examples above are necessarily non-polynomial under the
universal Euclidean criterion.
\end{proposition}

\begin{proof}
Let \(p=\max_i\deg h_i\), and choose \(i\) with \(\deg h_i=p\).  If
\(p>0\), the polynomial \(|\nabla h_i|^2\) has degree exactly \(2p-2\):
its leading term is the nonzero sum of squares of the derivatives of the
highest-degree homogeneous part of \(h_i\).  On the other hand, Theorem
\ref{thm:affine-gram} writes
\[
  |\nabla h_i|^2
  =
  c_i^0+\sum_{k=1}^d c_i^k h_k,
\]
whose degree is at most \(p\).  Therefore \(2p-2\leq p\), which implies
\(p\leq2\).
\end{proof}

\end{document}
